\documentclass[11pt]{article}
\usepackage[T1]{fontenc}
\usepackage{lmodern}
\usepackage[margin=1in]{geometry}
\usepackage{amsmath,amssymb,amsthm,booktabs}
\usepackage[colorlinks=true,urlcolor=blue,linkcolor=blue,citecolor=blue]{hyperref}
\usepackage{microtype}
\newtheorem{theorem}{Theorem}
\newtheorem{lemma}[theorem]{Lemma}
\newcommand{\Z}{\mathbb Z}
\newcommand{\Q}{\mathbb Q}
\newcommand{\OK}{\mathcal O_K}
\title{A direct split-prime bound for the tenth discriminant\\
\large Working refinement: \(q<10^6\)}
\author{Research note prepared with Codex}
\date{20 October 2026}
\begin{document}
\maketitle
\begin{abstract}
Assuming the determinant construction and Lemmas 4.2, 5.1, and 6.2 of
Banks, we obtain \(q<10^6\) for imaginary quadratic fields of class
number one. This improves the preceding cutoff \(5\cdot10^7\) by a
factor of 50. The new ingredient is a direct upper bound on the weighted
sum of split primes, obtained by counting primitive lattice points.
It replaces both the theta-integral zero construction and the
zero-dependent prime-bias estimate. Banks's determinant argument remains
an external input. No Heegner method or small-conductor zero computation
is used.
\end{abstract}

\paragraph{Status.}
This is a new working argument, with an exact rational certificate and
finite checks of the lattice-counting formulas. The independent check
of the preceding v2 argument does not cover this refinement. No new
independent mathematical review is recorded here.

\section{Result and inputs}

Let \(D<0\) be a fundamental discriminant and put
\[
 q=|D|,\qquad K=\Q(\sqrt D),\qquad R=\sqrt q,\qquad L=\log q.
\]
All logarithms are natural. Write \(\chi_D\) for the associated
quadratic character, so \(\chi_D(p)=1,-1,0\) according as a rational
prime splits, is inert, or ramifies in \(K\).

\begin{theorem}\label{main}
Assume Banks's determinant construction and Lemmas 4.2, 5.1, and 6.2
in \cite{banks}, including the interpolation results used in that
construction. If \(h(K)=1\), then
\[
                         q<1\,000\,000.
\]
\end{theorem}

We also use the classical explicit prime estimates of
Rosser--Schoenfeld \cite[Theorems 6 and 9]{rs}. In particular, for
\(X>H\ge319\),
\begin{equation}\label{mertens}
 \sum_{H<p\le X}\frac{\log p}{p}
 >\log X-\log H-\frac1{\log H},
 \qquad
 \vartheta(X):=\sum_{p\le X}\log p<2X.
\end{equation}
The first follows by subtracting the two bounds in their Theorem 6;
the constant term cancels. The second is weaker than their Theorem 9.

Unlike the previous argument \cite{v2}, this proof uses neither Banks's
Theorem 1.1 nor his Lemma 2.1 or Corollary 6.3. It uses the prime-power
divisibility in Lemma 6.2 directly. It does not independently establish
the cited determinant construction.

\section{Counting nonrational integers}

Choose the usual basis \(\OK=\Z+\Z\omega\), where
\(\operatorname{Re}\omega=u\in\{0,1/2\}\) and
\(\operatorname{Im}\omega=R/2\). Thus
\[
 Q(a,b):=|a+b\omega|^2=(a+ub)^2+\frac q4b^2.
\]
For \(x\ge0\), define
\[
 A(x)=\#\{(a,b)\in\Z\times\Z_{\ge1}:Q(a,b)\le x\}.
\]
The restriction \(b>0\) counts one half of the nonrational points.

\begin{lemma}\label{lattice}
For every \(x\ge0\),
\[
 \frac{\pi x}{R}-2\sqrt x-\frac{2\sqrt x}{R}
 \le A(x)\le
 \frac{\pi x}{R}+\frac{2\sqrt x}{R}.
\]
\end{lemma}
\begin{proof}
Put \(B=2\sqrt x/R\) and
\(f(t)=2\sqrt{x-qt^2/4}\) on \([0,B]\).
For each integer \(1\le b\le B\), the allowed real values of \(a\)
form a closed interval of length \(f(b)\). Its integer count lies
between \(f(b)-1\) and \(f(b)+1\). Since \(f\) is decreasing,
\[
 \int_0^B f(t)\,dt-f(0)
 \le\sum_{b=1}^{\lfloor B\rfloor}f(b)
 \le\int_0^B f(t)\,dt.
\]
The integral is \(\pi x/R\), and there are at most \(B\) intervals.
This gives both bounds, including when \(B<1\).
\end{proof}

Only primitive pairs can have prime norm. A small sieve already captures
enough of this restriction. Set
\[
 P=2\cdot3\cdot5\cdot7=210,\qquad
 A_P(x)=\#\{(a,b):b\ge1,\ Q(a,b)\le x,\ \gcd(a,b,P)=1\}.
\]

\begin{lemma}\label{primitive}
For every \(x\ge0\),
\[
 A_P(x)\le\frac{\pi\rho}{R}x+B_R\sqrt x,
 \qquad
 \rho=\frac{768}{1225},\quad
 B_R=\frac{88}{35}+\frac{192}{35R}.
\]
\end{lemma}
\begin{proof}
Homogeneity of \(Q\) and inclusion--exclusion give the exact identity
\[
 A_P(x)=\sum_{d\mid P}\mu(d)A(x/d^2).
\]
For positive coefficients use the upper bound of Lemma~\ref{lattice};
for negative coefficients use its lower bound. The area coefficient is
\[
 \sum_{d\mid P}\frac{\mu(d)}{d^2}
   =\prod_{p\mid P}(1-p^{-2})=\rho.
\]
Write
\[
 T=\sum_{d\mid P}\frac{|\mu(d)|}{d}=\frac{96}{35},\qquad
 U=\sum_{d\mid P}\frac{\mu(d)}d=\frac8{35}.
\]
The boundary coefficient is
\(2\sum_{\mu(d)=-1}d^{-1}+2T/R=T-U+2T/R=B_R\).
\end{proof}

\section{Direct bounds on split and ramified primes}

Assume \(h(K)=1\) and \(q\ge10^6\), and put \(a_0=q/4\).
If \(p\) splits, its two prime ideals of norm \(p\) are principal.
As the only units are \(\pm1\), there are exactly four generators
of norm \(p\). They give two distinct pairs with \(b>0\).
They are nonrational, and each satisfies \(\gcd(a,b,P)=1\), since
\(d\mid a,b\) would force \(d^2\mid Q(a,b)=p\).
Consequently,
\begin{equation}\label{split-count}
 S_+(X):=\sum_{\substack{p\le X\\\chi_D(p)=1}}\frac{\log p}{p}
 \le\frac12\sum_{\substack{b\ge1,\ Q(a,b)\le X\\\gcd(a,b,P)=1}}
                        \frac{\log Q(a,b)}{Q(a,b)}.
\end{equation}

All norms in this sum are at least \(a_0\). For \(X\ge a_0\),
partial summation, followed by Lemma~\ref{primitive}, therefore gives
\begin{align}
 S_+(X)
 &\le\frac12\left\{
 A_P(X)\frac{\log X}{X}
 +\int_{a_0}^X A_P(t)\frac{\log t-1}{t^2}\,dt\right\} \notag\\
 &\le\frac{\pi\rho}{4R}
       \bigl((\log X)^2-(\log a_0)^2+2\log a_0\bigr)
   +\frac{B_R(\log a_0+1)}{\sqrt{a_0}}.\label{split-precise}
\end{align}
In the last line we dropped the nonpositive boundary term
\(-B_R(\log X+2)/(2\sqrt X)\).
Here \(\log a_0>2\), and \(\log a_0+1<L\) because \(\log4>1\).
Also, for \(R\ge1000\),
\[
 \frac{\pi\rho}{4}<\frac{493}{1000},\qquad B_R<\frac{63}{25}.
\]
Thus the required direct split-prime bound is
\begin{equation}\label{split-final}
 S_+(X)<\frac{493}{1000R}(\log X)^2+\frac{126L}{25R}
             \qquad(X\ge q/4).
\end{equation}

A ramified prime likewise has a principal prime ideal of norm \(p\).
Its generator is nonrational, so \(p\ge q/4\). Hence
\begin{equation}\label{ramified}
 S_0(X):=\sum_{\substack{p\le X\\p\mid q}}\frac{\log p}{p}
 \le\frac4q\sum_{p\mid q}\log p\le\frac{4L}{q}.
\end{equation}
Neither this estimate nor the split-prime count requires a separate
reduction to prime discriminants.

Combining \eqref{mertens}, \eqref{split-final}, and \eqref{ramified},
for \(X>H\ge319\) and \(X\ge q/4\) we obtain
\begin{equation}\label{inert}
 \sum_{\substack{H<p\le X\\\chi_D(p)=-1}}\frac{\log p}{p}
 >\log X-\log H-\frac1{\log H}
  -\frac{493}{1000R}(\log X)^2
  -\left(\frac{126}{25R}+\frac4q\right)L.
\end{equation}

\section{The determinant comparison}

Take
\[
 H=10^9,\qquad N=q^{13},\qquad M=N^{4/3},\qquad
 z=\frac{52}{3},\qquad \log M=zL.
\]
For \(q\ge10^6\), Banks's hypotheses
\(q\ge9\), \(1\le H\le N\), and
\(H^{2/3}N^{4/3}\ge2\cdot97^2\) hold, as do \(M>H\) and \(M\ge q/4\).
Banks's determinant construction in \(E=K(\sqrt2)\) gives
\(\Delta\ne0\), the nonzero integer
\(J=\operatorname{N}_{E/\Q}(\Delta)\), and degree sums \(S_1,S_{23}\).
Lemma 4.2 gives
\[
 \frac{S_{23}}{S_1}\le\frac{768\cdot97^2}{10^9}
                  <r:=\frac{73}{10000},\qquad
 \frac{N^4M}{S_1}\le\frac{4\cdot97^2}{10^6}
                  <c:=\frac{19}{500}.
\]
Lemma 5.1 then implies
\begin{equation}\label{arch}
 \frac{\log|J|}{4S_1}
 \le(1+r)\left(\log N+\frac L2+\log8\right)
                         +\frac{2c\log N}{M}.
\end{equation}

For each prime \(p>H\) with \(\chi_D(p)=-1\), Lemma 6.2 supplies
\(p^{4E_p}\mid J\), where
\(E_p=\sum_{\alpha\in\mathcal A}\lfloor\alpha_1/p\rfloor
\ge S_1/p-N^4\). Multiplying the divisibilities for \(p\le M\), and
using \(\vartheta(M)<2M\), gives directly
\begin{equation}\label{padic}
 \frac{\log|J|}{4S_1}
 \ge\sum_{\substack{H<p\le M\\\chi_D(p)=-1}}\frac{\log p}{p}-2c.
\end{equation}
This step uses no zero-dependent corollary.

Compare \eqref{arch} with \eqref{padic} and insert \eqref{inert}
at \(X=M\). Since \(\log N=\tfrac34\log M\), division by \(L\)
gives the necessary inequality
\begin{align}
 \left(\frac14-\frac{3r}{4}\right)z
 &\le \frac{493}{1000}z^2\frac{L}{R}
       +\frac{1+r}{2}+\frac{126}{25R}+\frac4q \notag\\
 &\quad+\frac{\log H+(1+r)\log8+1/\log H+2c}{L}
       +\frac{26c}{M}.\label{necessary}
\end{align}

\section{A uniform rational contradiction}

The function \((\log q)/\sqrt q\) decreases for \(q>e^2\).
Using \(q\ge10^6\), we have the uniform estimates
\[
 \begin{gathered}
 \frac LR<\frac{13.816}{1000},\qquad L>13.815,\qquad R\ge1000,\qquad
 \frac{\log H}{L}\le\frac32,\\
 \log8<2.08,\qquad \frac1{\log H}<0.049,\qquad
 \frac{26c}{M}<10^{-6}.
 \end{gathered}
\]
The last inequality follows already from \(M\ge10^6\) and
\(26c=0.988<1\). The logarithmic estimates are certified by rational
Taylor bounds for the exponential function in the accompanying program.

The left side of \eqref{necessary} equals
\[
 \left(\frac14-\frac34\frac{73}{10000}\right)\frac{52}{3}
                  =\frac{127153}{30000}>4.238.
\]
Its right side is strictly less than the rational number
\begin{align*}
 \mathcal R={}&\frac{493}{1000}\left(\frac{52}{3}\right)^2
                       \frac{13.816}{1000}
             +\frac{1.0073}{2}+\frac32\\
 &+\frac{1.0073(2.08)+0.049+0.076}{13.815}
             +\frac{5.04}{1000}+\frac4{10^6}+\frac1{10^6}
             <4.216.
\end{align*}
Every terminating decimal here denotes its exact rational value.
The gap is greater than \(0.022\). This contradicts
\eqref{necessary}, proving Theorem~\ref{main}.

\clearpage
\section{Relation to earlier versions and verification}

The preceding v2 bound was \(q<5\cdot10^7\). It combined a theta-integral
zero with a zero-free estimate extracted from Banks. The present proof
instead obtains the necessary inert-prime mass directly from class
number one. Its small sieve removes lattice points whose two coordinates
share a factor among \(2,3,5,7\); it is not an exhaustive search through
discriminants.

The earlier finite-discriminant sieve covers \(q<10^{10}\), so its
unchanged certificates already cover the remaining interval below
\(10^6\). They leave \(3,4,7,8,11,19,43,67,163\). That finite
classification is separate from the analytic cutoff proved here.

All 26 exact arithmetic checks and 22 finite sanity cases pass.
To reproduce the exact rational certificate, run
\begin{verbatim}
python3 verify_direct_prime.py
\end{verbatim}
Add \texttt{--sanity} to enumerate lattice points in finite examples,
compare inclusion--exclusion with direct counting, and check that each
split prime has two generators with positive imaginary part in the
seven known fields of class number one with \(q>4\).

The finite checks do not prove the analytic lemmas or audit Banks's
interpolation argument. Codex generated the new derivation, this note,
and the checker. The preceding v2 cross-check remains scoped to v2.

\begin{thebibliography}{9}
\bibitem{banks}
W.~D.~Banks, manuscript \href{https://arxiv.org/abs/2610.15483v1}
{arXiv:2610.15483v1}, especially the determinant construction in
Sections 3--4 and Lemmas 4.2, 5.1, and 6.2.
\bibitem{rs}
J.~B.~Rosser and L.~Schoenfeld,
\emph{Approximate formulas for some functions of prime numbers},
Illinois J. Math. \textbf{6} (1962), 64--94, Theorems 6 and 9.
\href{https://doi.org/10.1215/ijm/1255631807}{doi:10.1215/ijm/1255631807}.
\bibitem{v2}
\emph{An explicit tenth-discriminant bound, version 2},
\href{https://github.com/twinheterodox/tenth-discriminant/releases/tag/v0.2.0}
{fixed release v0.2.0}, 20 October 2026.
\end{thebibliography}
\end{document}
